\subsection{Quadratic Jordan algebras and the composition equation}
\label{sec:quadratic-jordan-framework}

Quadratic-generated Laplacian algebras organize compositions of
arbitrary even degree. For a real symmetric matrix $A$, write
$q_A(x)=\langle Ax,x\rangle$.
Then
\begin{equation}\label{eq:quadratic-jordan-correspondence}
 \Delta q_A=2\tr A,\qquad
 \Gamma(q_A,q_B)=4q_{A\circ B},\qquad
 A\circ B=\tfrac12(AB+BA).
\end{equation}
Thus a quadratic-generated algebra $\cA$, containing $R=|x|^2$,
is Laplacian exactly when its quadratic part $J=\cA_2$ is a unital
Jordan algebra under $f\circ g=\Gamma(f,g)/4$
\cite[Proposition~37(b)]{MendesRadeschi}. The matrix realization matters: its multiplicities enter $\Delta q_A$.
The Jordan classification gives Clifford and real, complex, or
quaternionic Gram models; see
\cite[Theorems~B--C, Section~3 and Appendix~A]{MendesRadeschiQuadratic}.

The next proposition is the low-dimensional consequence of this
classification; it concerns algebras generated entirely in degree two.

\begin{proposition}[Three and four quadratic generators]
\label{prop:three-quadratic-normal-forms}
\label{prop:low-dimensional-quadratic-jordan}
Let $\cA$ be generated by homogeneous quadratics, contain $R$, and
be Laplacian. If $k=\dim\cA_2$ is three or four, the following
list is complete, up to orthogonal source coordinates and a linear
change of quadratic generators.
\begin{center}
\small
\begin{tabular}{cl>{\raggedright\arraybackslash}p{0.54\textwidth}}
\toprule
$k$ & $J$ & Quadratic coordinates \\
\midrule
3 & $\R^3$ & $|x_1|^2,|x_2|^2,|x_3|^2$ on three nonzero orthogonal blocks. \\
3 & $\operatorname{Sym}_2(\R)$ & $|x|^2,|y|^2,\langle x,y\rangle$, with $x,y\in\R^\ell$. \\
4 & $\R^4$ & Four nonzero orthogonal block radii. \\
4 & $\R\oplus\operatorname{Sym}_2(\R)$ & An independent block radius and the preceding real Gram coordinates. \\
4 & $\operatorname{Herm}_2(\mathbb C)$ & $|z|^2,|w|^2,\Re\langle z,w\rangle,\Im\langle z,w\rangle$, with $z,w\in\mathbb C^\ell$. \\
\bottomrule
\end{tabular}
\end{center}
Here $k$ is the number of minimal quadratic algebra generators;
algebraic relations may occur among them.
\end{proposition}

\begin{proof}
The correspondence \eqref{eq:quadratic-jordan-correspondence} embeds
$J$ as a unital Jordan subalgebra of the symmetric matrices. In the
formally real Jordan classification recalled in
\cite[Appendix~A]{MendesRadeschiQuadratic}, the simple factors of
real dimension at most four are $\R$, $\operatorname{Sym}_2(\R)$,
and $\operatorname{Herm}_2(\mathbb C)$, of dimensions $1,3,4$.
The spin factors of dimensions three and four are the latter two
matrix algebras. Taking direct sums gives precisely the displayed
list. A primitive one-dimensional factor acts by an orthogonal
projection; these projections split the source. The representation
argument in \cite[Section~3 and Appendix~A]{MendesRadeschiQuadratic}
gives repeated standard real or complex modules for the remaining
factors, hence their Gram coordinates. Finally, a degree-two element
of a quadratic-generated algebra is a linear combination of its
quadratic generators, proving the assertion about their minimal number.
\end{proof}

For the real two-column model use
$R=|x|^2+|y|^2$, $p=|x|^2-|y|^2$, and $q=2\langle x,y\rangle$.
Its image is $\{R\geq0:p^2+q^2\leq R^2\}$ for $\ell\geq2$;
for $\ell=1$ only the boundary $R^2=p^2+q^2$ is attained.
The complex two-column model similarly has coordinates
$R,p,q_1,q_2$, with $q_1+iq_2=2\langle z,w\rangle$ and
$p=|z|^2-|w|^2$. Its image is
$\{R\geq0:p^2+q_1^2+q_2^2\leq R^2\}$ for $\ell\geq2$;
for $\ell=1$ only its boundary is attained. These statements follow by realizing positive
semidefinite $2\times2$ Gram matrices of the allowed rank.

An arbitrary quadratic span need not be closed: on $\R^3$,
$\operatorname{span}\{R,x_1^2,x_1x_2\}$ omits
$\Gamma(x_1x_2,x_1x_2)=x_1^2+x_2^2$. Adjoining $x_2^2$ gives
$\R\oplus\operatorname{Sym}_2(\R)$. Every finite list has a unital
matrix Jordan closure of dimension at most $N(N+1)/2$; restricting
its size is essential, since all even-degree homogeneous polynomials
are polynomials in the full collection $x_ix_j$.

Choose a basis $z=(z_1,\ldots,z_k)$ of $J$ and put
\begin{equation}\label{eq:quadratic-jordan-profile-data}
 \begin{gathered}
 B_{ij}(z)=\tfrac14\Gamma(z_i,z_j),\qquad
 b_i=\tfrac14\Delta z_i,\\
 A_P=\sum_{i,j}B_{ij}P_iP_j,\qquad
 \mathscr D P=\sum_{i,j}B_{ij}P_{ij}+\sum_i b_iP_i.
 \end{gathered}
\end{equation}
The entries of $B$ are linear and the $b_i$ are constant.
The quotient calculus in \cref{thm:finite-reduction} gives
\begin{equation}\label{eq:quadratic-jordan-profile-pullback}
 |\nabla P(z)|^2=4A_P(z),\qquad
 \Delta P(z)=4(\mathscr D P)(z),\qquad
 \cL(P(z))=8\mathscr T(P)(z),
\end{equation}
where
\begin{equation}\label{eq:three-quadratic-profile-numerator}
 \mathscr T(P)=2A_P\mathscr D P
       -\sum_{i,j}B_{ij}P_i\partial_jA_P.
\end{equation}
These identities are evaluated on the actual generator image, also
when the generators satisfy relations.

\begin{proposition}[Composition criterion]
\label{prop:three-quadratic-profile-criterion}
\label{prop:quadratic-jordan-profile-criterion}
Let $P$ be a real homogeneous polynomial of degree $d$ and set
$F=P(z)$. Its regular zero set is minimal if and only if
$\mathscr T(P)=0$ on the part of the generator image where
$P=0$ and $A_P>0$. In particular, $P\mid\mathscr T(P)$ is
sufficient. It is also necessary if $P$ is irreducible and has a
regular zero in an open subset of the image where $B$ is positive
definite. In that case the quotient has degree $2d-2$, while $F$
has degree $2d$.
\end{proposition}

\begin{proof}
The first assertion follows from
\eqref{eq:quadratic-jordan-profile-pullback} and the level-set
mean-curvature identity. Positive definiteness makes the generator
map a submersion and gives $A_P>0$ at regular profile zeros.
Minimality therefore makes $\mathscr T(P)$ vanish on a relatively
open smooth real patch of $P=0$. The patch-divisibility argument in
\cref{prop:minimizing-irreducibility} gives $P\mid\mathscr T(P)$.
Its degree is $3d-2$, which proves the last assertion.
\end{proof}

For orthogonal radii $u_i=|x_i|^2$, with $n_i=\dim V_i$, the data are
\begin{equation}\label{eq:three-radii-profile-operator}
 B=\operatorname{diag}(u_i),\qquad
 \mathscr D P=\sum_i u_iP_{ii}+\tfrac12\sum_i n_iP_i.
\end{equation}
Thus the criterion specializes to
\begin{equation}\label{eq:three-radii-minimality}
 P\ \bigm|\
 \left(\sum_i u_iP_i^2\right)
 \left(2\sum_i u_iP_{ii}+\sum_i n_iP_i\right)
 -\sum_i u_iP_i\partial_i\left(\sum_j u_jP_j^2\right),
\end{equation}
under the irreducibility and interior-zero hypotheses above.

\subsubsection{Clifford profiles with a radial variable}

For a symmetric Clifford system on $\R^N$ with $s\geq2$, let
$p_i(x)=\langle P_ix,x\rangle$ for $1\leq i\leq s$.
The algebra $\R[R,p_1,\ldots,p_s]$ has
\begin{equation}\label{eq:three-spin-profile-operator}
 B=\begin{pmatrix}R&p^T\\p&R\Id_s\end{pmatrix},\qquad
 \mathscr D P=R(P_{RR}+\Delta_pP)
      +2\sum_i p_iP_{Rp_i}+\tfrac N2P_R.
\end{equation}
This follows from \eqref{eq:Clifford-pullback-identities}; the real
and complex two-column models have $(s,N)=(2,2\ell)$ and $(3,4\ell)$.

If a homogeneous profile $f(p)$ is independent of $R$, the identity
\eqref{eq:pullback-defect} reads
\begin{equation}\label{eq:quadratic-Clifford-profile-transfer}
 \cL_x(f(p))=16R^2(\cL_pf)(p)
             -8d f(p)|\nabla_pf|^2(p),\qquad d=\deg f.
\end{equation}
Thus a degree-$d$ minimal target cone yields a degree-$2d$ minimal
source equation, and a harmonic target gives a harmonic source.
This is the Clifford specialization of Riedler's cone-pullback
theorem \cite[Theorem~5.1.4]{RiedlerThesis25}; its subcalibration
transfer is \cref{thm:clifford-area-transfer}. In particular, minimal
cubic targets supply sextic equations. Effective factors and the
variational conclusions are determined by the corresponding results
in \cref{sec:quartic-application}.

Allowing the radial coordinate also includes the OT--FKM splitting
in \eqref{bh:centered}. For a homogeneous profile $P(t,y)$ of degree
$d$, define
\[
 K_P=|\nabla_yP|^2-P_t^2,\qquad
 \square P=\Delta_yP-P_{tt},\qquad
 \Gamma_L(P,K_P)=\nabla_yP\cdot\nabla_yK_P-P_t(K_P)_t.
\]
On $F=P(R,p)=0$, the same profile table gives
\begin{align}
 |\nabla F|^2&=4R K_P,\notag\\
 \cL(F)&=16R^2\left(K_P\square P
                  -\tfrac12\Gamma_L(P,K_P)\right)
             +8(N-2)RK_PP_t.
 \label{eq:radial-Clifford-profile-equation}
\end{align}
For clarity, Euler's identity first gives, on the whole source,
\[
 |\nabla F|^2=4RK_P+8dPP_t,\qquad
 \Delta F=4R\square P+(2N+8d-8)P_t.
\]
Differentiate the first expression before setting $P=0$;
$\Gamma(F,R)=4dF$ and the homogeneity of $K_P$, of degree $2d-2$,
then give \eqref{eq:radial-Clifford-profile-equation}.
For $P=|y|^2-ct^2$, $0<c<1$, this equation selects
$c=2(s-1)/(N-2)$, the familiar minimal isoparametric level.
The splitting quartic instead has profile
$\tfrac14(\sum_{i=1}^m p_i^2+cp_0^2-cR^2)$, while Riedler's
quartic is independent of $R$. Both therefore belong to the same
composition calculus.

\begin{remark}[An elementary radial profile]
\label{sec:three-quadratic-profiles}
In the real two-column model, the radial variable gives quartic minimal levels: for
$\ell>2$ set
\[
 W=p^2+q^2,\qquad c=\frac1{\ell-1},\qquad
 F=W-cR^2,\qquad \Phi=R^2-2W.
\]
Equation \eqref{eq:radial-Clifford-profile-equation} selects this $c$;
$|\nabla\Phi|^2=16R^3$ and $\Delta\Phi=8(\ell-3)R$ identify
$F=0$ as its minimal isoparametric level. Its smallest algebra is
$\R[R,F]$, since $W=F+cR^2$. This minimal level is covered by
\cref{sec:isoparametric-reduction}.
\end{remark}


\subsubsection{Real and complex Gram profiles}

Let $X=[x_1\ \cdots\ x_r]\in\R^{p\times r}$ and set
$h_{ij}=\langle x_i,x_j\rangle$. The independent entries of
$H=X^TX$ span $\operatorname{Sym}_r(\R)$, with $R=\tr H$, and
\begin{equation}\label{eq:real-Gram-differential-table}
 \begin{split}
 \Delta h_{ij}&=2p\delta_{ij},\\
 \Gamma(h_{ij},h_{kl})
 &=\delta_{ik}h_{jl}+\delta_{il}h_{jk}
   +\delta_{jk}h_{il}+\delta_{jl}h_{ik}.
 \end{split}
\end{equation}
Indeed, $\nabla_{x_a}h_{ij}=\delta_{ai}x_j+\delta_{aj}x_i$;
taking inner products and traces proves the table, including the
diagonal entries. For $Z\in\mathbb C^{p\times r}$, use the real
Euclidean metric and write
$h_{ij}=\sum_a\overline{z_{ai}}z_{aj}$, so $H=Z^*Z$.
The complex-bilinear extension of $\Gamma$ satisfies
\begin{equation}\label{eq:complex-Gram-differential-table}
 \Delta h_{ij}=4p\delta_{ij},\qquad
 \Gamma(h_{ij},h_{kl})
       =2(\delta_{jk}h_{il}+\delta_{il}h_{kj}).
\end{equation}
To verify it, use $\Delta=4\sum_{a,b}\partial_{z_{ab}}
\partial_{\bar z_{ab}}$ and
$\Gamma(f,g)=2\sum_{a,b}(f_{z_{ab}}g_{\bar z_{ab}}
+f_{\bar z_{ab}}g_{z_{ab}})$. Taking real and imaginary parts
gives the real table on $\operatorname{Herm}_r(\mathbb C)$.
These are the standard Gram realizations in
\cite[Section~3 and Section~7]{MendesRadeschiQuadratic}.

Their generator images consist of positive semidefinite Gram matrices
of rank at most $p$. For $p\geq r$ the positive definite part is
open and the Gram map is a submersion there: for a prescribed
symmetric or Hermitian increment $A$, the variation
$\dot X=\tfrac12 XH^{-1}A$, or
$\dot Z=\tfrac12 ZH^{-1}A$, has $\dot H=A$.
Consequently \cref{prop:quadratic-jordan-profile-criterion} applies
directly to profiles with a regular zero in this positive definite
part. For $p<r$ the same differential identities apply on the
rank-constrained image. The constructions below select minimal
profiles from these differential tables.
